% Source: https://artofproblemsolving.com/community/c3h1404901_even_more_lithuanian_math_kangaroo_problems__junior_2010
% Stable identity: aops-topic-1404901-q29
Put the centre of the rectangle at the origin, so its sides are $x=\pm4$ and $y=\pm2$. The right arc touches $x=4$ at $(4,0)$ and has radius $1$, so its centre is $O_1=(3,0)$. The top arc touches $y=2$ at $(0,2)$ and has radius $R$, so its centre is $O_2=(0,2-R)$ (below the oval's centre, since $R>2$).

\begin{asy}
usepackage("amsmath");
size(6.5cm);
defaultpen(fontsize(10pt));
real R=6;
draw((-4,-2)--(4,-2)--(4,2)--(-4,2)--cycle,gray+dashed);
real ang=degrees(atan2(4,3));
draw(arc((3,0),1,-ang,ang),blue+linewidth(0.9)); draw(arc((-3,0),1,180-ang,180+ang),blue+linewidth(0.9));
draw(arc((0,2-R),R,ang,180-ang),red+linewidth(0.9)); draw(arc((0,R-2),R,180+ang,360-ang),red+linewidth(0.9));
draw((0,2-R)--(0,2),red+dashed); draw((0,2-R)--(3.6,0.8),red+dashed); draw((3,0)--(4,0),blue+dashed);
dot((3,0),blue); dot((0,2-R),red); dot((3.6,0.8));
label("$O_1$",(3,0),dir(290),blue); label("$O_2$",(0,2-R),E,red); label("$T$",(3.6,0.8),NE);
label("$1$",(3.72,0),N,blue); label("$R$",(0,-1),W,red);
label("$8$",(-2,-2),S,gray); label("$4$",(-4,0),W,gray);
label("\mbox{}",(min(currentpicture,user=true)+max(currentpicture,user=true))/2,black); // keeps the text after the figure black when LaTeX typesets the labels inline
\end{asy}

The two arcs are tangent at a point $T$, and the small circle lies inside the big one, so the centres and $T$ are collinear and $O_1O_2=R-1$:
\[
3^2+(R-2)^2=(R-1)^2\quad\Longrightarrow\quad 9+R^2-4R+4=R^2-2R+1\quad\Longrightarrow\quad R=6 .
\]
So the radius of the big arcs is $\boxed{6}$, choice (A). (Then $O_2=(0,-4)$ and $T=(3.6,0.8)$.)
